\section{Introduction}
A finite command catalog creates a distinction between the resources allocated to a decision and the number of commands used to implement it. A provider may deliver different expected obligations to different customer histories while installing only a small common set of terminal commands. We call the installed set a command book. Each installed command incurs a fixed charge. Individual expected caps and ceilings on every realized obligation restrict how an accepted aggregate promise can be implemented. Choosing expected targets first and a catalog afterward can therefore miss both a discrete cost and a feasibility constraint.

This paper establishes a tariff-sensitive complexity frontier for joint resource allocation and implementation. The main contrast holds even when rewards are linear and preliminary service is free. With arbitrary command-specific opening fees, exact optimization remains parameterized-hard in both the command allowance and the number of distinct expected caps. With a standard fee and only a few exceptional commands, however, we obtain an exact fixed-parameter algorithm whose polynomial exponent does not depend on the number of exceptions. Uniform fees give a polynomial-time algorithm with no discretization of capacities or the accepted promise. Thus a small command allowance does not by itself make the problem tractable, whereas a simple installation tariff can do so without restricting the physical heterogeneity.

The positive result follows from a capacity representation inside the original model. A selected book's gross value depends only on its attainable aggregate terminal capacity. This capacity telescopes along consecutive selected commands. Conditional on the installed exceptional commands and the book size, the total opening charge is fixed; it is then sufficient to retain one maximum-capacity label at each state. A recovered label is not merely a numerical value: it supplies terminal lotteries and preliminary service that meet the original promise and every individual constraint. The same tables compute the provider's exact decision frontier as a uniform installation fee changes. The dynamic-programming principle is familiar; its applicability here follows from the model-specific capacity and fixed-charge identities.

The exact selected-boundary representation also supports the general concave model. We characterize capacity conservation on acyclic graphs and extend the recognition criterion to paths with an explicit edge allowance. Complementing resources merges feasible starting points into a single deficit state. Rounding and repair preserve the selected path, its command count and the accepted promise. These results retain the original-feasible additive guarantee even where the tariff algorithm is inapplicable. A corridor-rehabilitation model illustrates the conservation and repair mechanism outside renewal commands, with its physical quantities and separability assumptions stated explicitly in the companion.

Price-based certificates and resource approximation play different roles. We characterize exactly when one price-active book supports the accepted promise, and bound the residual gap by the distance to that book's maximizing target interval. This explains why price methods can be strong on root-supported instances without implying that their general search trees are small. A uniform positive installation fee can still create a price gap even though the tariff algorithm solves the problem exactly. The earlier deficit implementation's unfavorable timing results remain part of the evidence; an accuracy guarantee does not establish practical speed superiority.

The computational study evaluates these distinctions rather than claim operational adoption. It preserves the earlier records, distinguishes distinct model inputs from tolerance-tagged cases, and adds prospectively frozen standard-tariff, exception-tariff and curved-return families under three end-to-end allowances. The direct mixed-integer convex formulation has no service-envelope approximation. Numerical bounds from its solver remain separate from independently checked rational certificates. Setup, serialization and verification consume the same allowance as optimization; unsuccessful requests remain in the denominator.

\subsection{Relationship to prior work}
Lagrangian constrained-path bounds, label search and resource-state approximation are established tools \citep{HandlerZang1980,BeasleyChristofides1989,Fisher1981,Hassin1992,LorenzRaz2001,PuglieseGuerriero2013}. Contemporary bucket-graph labeling organizes dominance and reduced-cost elimination for routing subproblems \citep{Sadykov2021}; heuristic-guided constrained search addresses large networks and multiple resource limits \citep{Ahmadi2025}. Those methods are not displaced by our deficit dynamic program. Our graph edges instead carry allocatable continuous resources, and our claims concern an exact representation of individual obligations, a recognizable conservation restriction, and recovery of an original equality. Generic path algorithms do not establish those model identities. No timing comparison against a routing implementation is inferred from our allocation experiments.

Continuous allocation on a fixed book belongs to separable resource allocation \citep{Patriksson2008}. Fixed-charge nonlinear allocation also admits perspective reformulations and branch-and-Benders-cut approaches \citep{Yang2025}. These are important formulation and algorithmic antecedents, not results that opening fees generally become easy. Our positive theorem exploits the additional ordered-capacity identity and the fact that conditioning on exceptional commands makes a given book size have a common charge. Its hardness counterpart concerns parameterized complexity, rather than the presence of binary variables alone. Parameterized knapsack classifications distinguish structural parameters, numerical values and their encodings \citep{Gurski2019}; their bounds do not transfer to our shared-book constraints without a reduction. We supply that reduction directly inside the original policy model.

Small distributional support is not a small shared installation catalog. Moment-set extreme-point results bound individual supports \citep{Winkler1988}; they do not charge for their union. Auction menu complexity additionally imposes incentive constraints absent here \citep{HartNisan2013,Babaioff2016}. Choice-based assortment optimization changes demand through the offered set \citep{Rusmevichientong2010}, whereas our history probabilities are fixed and terminal probabilities are controlled. Our reward perturbation certificate holds feasibility inputs fixed; it is not a statistical distributional-robustness claim \citep{EsfahaniKuhn2018}.

Finally, concave convolution uses classical matrix-searching structure \citep{Aggarwal1987}; the finite-support proof permits holes created by merged sources. General medoid selection and one-dimensional weighted-median clustering are antecedent algorithms, not independent novelty claims \citep{Arya2004,Wu1991}. Their role is to make the model's oscillation-based grouping bound operational. Table~\ref{tab:frontier60} separates these antecedents from the results established here. The paper's scientific contribution is the exact representation and tariff-sensitive frontier, with conservation, approximation and certification supplying implementable consequences.
